\subsection{A Complete Replay Decision}
\label{app:prompt-trace}
The following first-iteration repair concerns a modular-arithmetic skill. The current task is ``Is $385^{1980}+18^{1980}$ a perfect square?'' and the skill is annotated as required. This example distinguishes the predicate immediately before the repair from the cold-start baseline used by the held-out gate.

\paragraph{Predicate before repair.}
\begingroup
\small\ttfamily\raggedright
The task asks for a lower bound on the count of variables divisible by a specific modulus m for integer solutions of an equation.
\par
\endgroup

\paragraph{Generated alternatives.}
\noindent\textbf{Candidate 1.}
\begingroup
\small\ttfamily\raggedright
The task asks whether a specific integer expression is a perfect square, requiring verification that the value is NOT equal to any non-zero square or identification of an impossibility.
\par
\endgroup
\noindent\textbf{Candidate 2.}
\begingroup
\small\ttfamily\raggedright
The task asks to determine possible residue classes modulo an integer to prove an expression cannot be a perfect square or to enumerate remainders of an exponential expression.
\par
\endgroup
\noindent\textbf{Candidate 3.}
\begingroup
\small\ttfamily\raggedright
The task asks to find all integer solutions to an equation in the form of checking if an expression can be a perfect square by analyzing divisibility and residue constraints.
\par
\endgroup

\paragraph{Replay evidence.}
The five relevant tasks ask for: (i) all values of $4^n+15n-1\bmod9$ for integer $n\geq0$; (ii) the smallest positive integer whose first digit can be moved to the end to multiply its value by $1.5$; (iii) integer solutions of $19x^3-84y^2=1984$; (iv) the current perfect-square question; and (v) $3^{105}+4^{105}\bmod13$.

The five irrelevant tasks ask for: (i) integer $x$ for which $x$ and $8x^2+1$ are prime; (ii) whether $n=x^2+y^2$ implies that $2n$ is a sum of two squares; (iii) integer $m$ for which $1000^m-1$ divides $1978^m-1$; (iv) possible $a+b\bmod d^2$ when $d=\gcd(a,b)$ and $(a+1)/b+(b+1)/a$ is an integer; and (v) whether $(1/2)(3/4)\cdots(99/100)>0.1$. These are condensed task descriptions; the current question and predicate strings above are exact.

\begin{table}[t]
\centering\small
\begin{tabular}{lrrrr}
\toprule
Predicate & TP & FP & Errors & Risk\\
\midrule
Before & 0 & 0 & 5 & 0.50\\
Candidate 1 & 1 & 0 & 4 & 0.40\\
Candidate 2 & 3 & 0 & 2 & 0.20\\
Candidate 3 & 2 & 0 & 3 & 0.30\\
\bottomrule
\end{tabular}
\caption{One repair on five positive and five negative replay tasks. In this example, all three candidates correct the current failure and preserve already-correct replay judgments.}
\label{tab:appendix-repair-trace}
\end{table}

The old predicate is false on all ten tasks. Candidate 2 recovers three positives, including the current task, without introducing a false positive. Its replay risk is lowest, so it replaces the old predicate. Two positive tasks remain misclassified; the repair does not require perfect replay accuracy.

\paragraph{Held-out retention.}
The cold-start predicate for this skill was:
\begingroup
\small\ttfamily\raggedright
The task requires determining whether integer solutions exist, enumerating possible residue values, proving universal divisibility, or finding terms in a periodic sequence based on modular constraints.
\par
\endgroup

At the end of the first training pass, Candidate 2 is still the skill's candidate predicate. On the same 44 held-out tasks retrieving this skill, correctness increases from $33/44$ for the cold-start predicate to $39/44$ for the candidate. The strict gate therefore retains it. This comparison is the iteration-end skill decision, not an extra gate immediately after each individual repair.
